\chapter{Estimation par maximum de vraisemblance}
\label{ch:likelihood}

\begin{lecture}
L'apprentissage compare les probabilités annoncées aux réponses observées. Un
paramétrage reçoit une vraisemblance élevée lorsqu'il attribue une forte
probabilité aux étiquettes effectivement présentes dans le fichier. La
log-perte est la même comparaison écrite sous la forme d'un coût à minimiser.
\end{lecture}

On suppose les observations indépendantes conditionnellement à leurs variables
explicatives et
\[
  Y_i\mid X_i=\x_i\sim
  \operatorname{Bernoulli}(p_i),\qquad p_i=\sigma(\thv^\top\xt_i).
\]
Ce modèle de Bernoulli conditionnel et son estimation par maximum de
vraisemblance sont la construction standard de la régression logistique
\cite[sec.~4.4.1]{hastie2009,cox1958}.

À une valeur donnée de $\thv$ correspondent donc trois objets successifs : le
vecteur des scores $\X\thv$, le vecteur des probabilités $\mathbf p$ et un coût scalaire
$J(\thv)$ qui résume l'accord avec toutes les cibles. L'apprentissage réévalue
ces objets après chaque modification de $\thv$ ; les observations et leurs
cibles restent fixes.

La masse de probabilité d'une Bernoulli peut s'écrire sous une forme unique pour
les deux valeurs possibles de $y_i$ :
\[
  \Pp(Y_i=y_i\mid X_i=\x_i;\thv)
  =p_i^{y_i}(1-p_i)^{1-y_i}.
\]
Si $y_i=1$, cette expression devient $p_i$ ; si $y_i=0$, elle devient $1-p_i$.
Sous l'hypothèse d'indépendance conditionnelle des observations, la probabilité
jointe des réponses observées est le produit de ces contributions. C'est ce
passage --- et non une analogie avec les moindres carrés --- qui produit la
vraisemblance suivante.

Cette indépendance concerne les élèves une fois leurs caractéristiques fixées.
Elle permet d'écrire, pour deux observations,
\begin{align*}
  &\Pp(Y_1=y_1,Y_2=y_2\mid X_1=\x_1,X_2=\x_2;\thv)\\
  &\qquad=\Pp(Y_1=y_1\mid X_1=\x_1;\thv)
   \Pp(Y_2=y_2\mid X_2=\x_2;\thv).
\end{align*}
Elle serait douteuse si plusieurs lignes correspondaient au même élève ou à des
jumeaux dont les réponses restent liées après prise en compte des variables. Le
produit de vraisemblance n'est donc pas une identité universelle : il traduit une
hypothèse sur la manière dont les données ont été recueillies.
La vraisemblance et sa log-vraisemblance sont
\begin{align*}
  L(\thv)&=\prod_{i=1}^m p_i^{y_i}(1-p_i)^{1-y_i},\\
  \log L(\thv)&=\sum_{i=1}^m
  \left[y_i\log p_i+(1-y_i)\log(1-p_i)\right].
\end{align*}
Maximiser $L$ équivaut à minimiser la log-perte moyenne
\[
  J(\thv)=-\frac1m\log L(\thv)
  =-\frac1m\sum_{i=1}^m\left[y_i\log p_i+
  (1-y_i)\log(1-p_i)\right].
  \tag{1}\label{eq:loss}
\]
L'équation~\eqref{eq:loss} est l'opposé de la log-vraisemblance binomiale,
également appelée log-perte ou entropie croisée binaire
\cite[équations~4.19--4.20]{hastie2009}.
Cette quantité est l'opposé du logarithme de la probabilité attribuée par le
modèle aux étiquettes observées.

\begin{table}[htbp]
\centering
\begin{tabular}{ccccc}
\toprule
$i$ & $y_i$ & $p_i$ & probabilité attribuée au résultat observé & $\ell_i$ \\
\midrule
$1$ & $1$ & $0{,}90$ & $0{,}90$ & $-\log(0{,}90)=0{,}105$ \\
$2$ & $0$ & $0{,}20$ & $1-0{,}20=0{,}80$ & $-\log(0{,}80)=0{,}223$ \\
$3$ & $1$ & $0{,}10$ & $0{,}10$ & $-\log(0{,}10)=2{,}303$ \\
\midrule
\multicolumn{4}{r}{perte moyenne $J$} & $0{,}877$ \\
\bottomrule
\end{tabular}
\caption{Calcul de la log-perte sur trois observations. La troisième prédiction,
fausse et assurée, domine la moyenne ; c'est le comportement recherché par le
maximum de vraisemblance.}
\label{tab:loss-example}
\end{table}

\begin{figure}[htbp]
  \centering
  \begin{tikzpicture}
    \begin{axis}[courseplot,axis lines=left,
      xlabel={probabilité prédite $p$},ylabel={perte $\ell(y,p)$},
      xmin=0,xmax=1,ymin=0,ymax=5,grid=major,grid style={gridgray!50},
      samples=250,domain=.007:.993,legend pos=north west]
      \addplot[accent,very thick]{-ln(x)}; \addlegendentry{$y=1$: $-\log p$}
      \addplot[warning,very thick]{-ln(1-x)}; \addlegendentry{$y=0$: $-\log(1-p)$}
    \end{axis}
  \end{tikzpicture}
  \caption{Une prédiction correcte et assurée coûte presque zéro ; une
  prédiction erronée et assurée reçoit une pénalité non bornée.}
  \label{fig:loss}
\end{figure}

\begin{figure}[htbp]
\centering
\begin{tikzpicture}
\begin{axis}[courseplot,axis lines=left,
  xlabel={probabilité prédite $p$ pour une observation $y=1$},ylabel={pénalité},
  xmin=.02,xmax=1,ymin=0,ymax=4.2,grid=major,grid style={gridgray!45},
  legend pos=north east,samples=240,domain=.02:1]
  \addplot[warning,very thick]{-ln(x)};
  \addlegendentry{log-perte $-\log p$}
  \addplot[accent,dashed,very thick]{(1-x)^2};
  \addlegendentry{erreur quadratique $(1-p)^2$}
\end{axis}
\end{tikzpicture}
\caption{Comparaison pour $y=1$. Les deux pertes sont minimales en $p=1$, mais
la log-perte pénalise beaucoup plus fortement une probabilité presque nulle
attribuée à l'événement réellement observé. Elle découle en outre directement
de la vraisemblance de Bernoulli.}
\label{fig:loss-comparison}
\end{figure}

\begin{procedure}{Évaluer la log-perte}
\textbf{Entrées :} $m$ probabilités $p_i$ et $m$ cibles $y_i\in\{0,1\}$.
\textbf{Sortie :} un nombre $J(\thv)$ ; une valeur plus faible indique un meilleur
accord avec les étiquettes observées.
\begin{enumerate}
  \item calculer les scores $z_i$, puis les probabilités $p_i$ ;
  \item pour chaque ligne, prendre $a_i=p_i$ si $y_i=1$, sinon $a_i=1-p_i$ ;
  \item calculer $\ell_i=-\log(a_i)$ ; par exemple, $a_i=0{,}9$ donne
        $\ell_i\approx0{,}105$ et $a_i=0{,}1$ donne $\ell_i\approx2{,}303$ ;
  \item calculer $J=m^{-1}\sum_i\ell_i$ et vérifier qu'il reste fini.
\end{enumerate}
\end{procedure}

\section{Gradient et correction des coefficients}

\begin{lecture}
Le gradient indique comment une petite modification de chaque coefficient ferait
varier le coût. Chaque observation contribue selon son erreur $p_i-y_i$ et selon
ses caractéristiques. La somme de ces contributions fournit la correction de
tous les coefficients en une opération matricielle.
\end{lecture}

Pour chaque observation, le programme calcule d'abord l'écart $p_i-y_i$ entre
la probabilité produite et la cible binaire. Il pondère ensuite cet écart par les
caractéristiques de l'observation, puis additionne les contributions. Le calcul
complet s'écrit
\[
  \boxed{\quad\nabla J(\thv)=\frac1m\X^\top
  \bigl(\sigma(\X\thv)-\y\bigr).\quad}
  \tag{2}\label{eq:gradient}
\]
Cette expression est le score de vraisemblance sous forme matricielle
\cite[équation~4.21]{hastie2009}.
Sa dérivation par la règle de chaîne figure à
l'annexe~\ref{app:gradient-derivation}.
Le gradient a bien la dimension $n+1$ : $\X^\top$ est de taille
$(n+1)\times m$ et le vecteur des résidus probabilistes est de taille $m$.

Le terme $(p_i-y_i)\xt_i$ permet une lecture locale. Pour un positif
$y_i=1$ sous-estimé à $p_i=0{,}1$, le résidu vaut $-0{,}9$ : l'observation exerce
une correction importante. Pour un positif déjà estimé à $0{,}99$, le résidu
vaut seulement $-0{,}01$. Le gradient agrège ces corrections en tenant compte de
la direction $\xt_i$ de chaque observation.

\begin{procedure}{Calculer le gradient}
\textbf{Entrées :} $\X\in\R^{m\times(n+1)}$, $\y\in\{0,1\}^m$ et
$\thv\in\R^{n+1}$. \textbf{Sortie :} $\mathbf g\in\R^{n+1}$, dont chaque
composante indique la correction associée à un coefficient.
\begin{enumerate}
  \item calculer $\mathbf z=\X\thv$, de longueur $m$, puis
        $\mathbf p=\sigma(\mathbf z)$ ;
  \item calculer $\mathbf r=\mathbf p-\y$ ; un positif prédit à $0{,}1$ contribue
        par un écart $-0{,}9$ ;
  \item calculer $\mathbf g=\X^\top\mathbf r/m$ ;
  \item vérifier que $\mathbf g$ et $\thv$ ont tous deux $n+1$ composantes ;
  \item avant l'entraînement complet, comparer quelques composantes à la
        différence finie de la figure~\ref{fig:finite-difference}.
\end{enumerate}
\end{procedure}

\section{Convexité, unicité et séparabilité}

\begin{lecture}
Une fonction convexe a la géométrie générale d'une cuvette : descendre localement
ne conduit pas vers un faux minimum isolé. Une cuvette peut toutefois posséder un
fond plat, donc plusieurs paramètres équivalents. Avec des classes parfaitement
séparées, le minimum non pénalisé peut même être repoussé à l'infini.
\end{lecture}

En notant $W=\operatorname{diag}(p_i(1-p_i))$, la Hessienne vaut
\[
  \nabla^2J(\thv)=\frac1m\X^\top W\X.
\]
Cette Hessienne est la matrice d'information observée, à un facteur de
normalisation près \cite[équations~4.22--4.25]{hastie2009}.
Les coefficients diagonaux $p_i(1-p_i)$ sont positifs ou nuls ; la Hessienne est
donc positive semi-définie et $J$ est convexe. Tout minimum local est global. La
preuve matricielle est donnée à l'annexe~\ref{app:convexity-proof}. La stricte
convexité, et donc l'unicité du minimiseur, exige toutefois des conditions de rang
et l'absence de directions dégénérées.

La distinction entre convexité et stricte convexité est importante. Une fonction
convexe ne possède pas de minimum local trompeur, mais elle peut présenter un
fond plat contenant plusieurs minimiseurs. Si deux colonnes de $\X$ sont
identiques, par exemple, certains changements opposés de leurs coefficients
laissent les scores inchangés. La prédiction peut alors être unique tandis que
le vecteur de paramètres ne l'est pas.

En une dimension, $J(\theta)=(\theta-1)^2$ est strictement convexe : sa courbure
$J''(\theta)=2$ est positive et son unique minimum est $\theta=1$. La fonction
$J(\theta_1,\theta_2)=(\theta_1+\theta_2-1)^2$ est seulement convexe : tous les
couples vérifiant $\theta_1+\theta_2=1$ minimisent le coût. Cette seconde
situation reproduit ce qui se passe lorsque deux colonnes de la matrice de design
apportent exactement la même information.

\begin{vigilance}[Séparation complète]
Si un hyperplan sépare parfaitement les deux classes, la vraisemblance non
pénalisée peut croître lorsque $\|\thv\|\to\infty$ sans atteindre de maximiseur
fini. La convexité seule n'assure donc pas l'existence d'un minimiseur fini. Une
pénalisation $L^2$, ou une autre méthode adaptée, rétablit en
général un problème bien posé \cite{albert1984}.
\end{vigilance}

Exemple : sur une caractéristique $x$, supposons que tous les négatifs vérifient
$x<0$ et tous les positifs $x>0$. Le score $z=10x$ classe déjà les observations
avec assurance ; $z=100x$ rend les probabilités encore plus proches de $0$ ou de
$1$ et augmente encore la vraisemblance. Aucun coefficient fini n'est désigné
comme optimum : le modèle pousse le coefficient vers $+\infty$. La pénalisation
ajoute un coût croissant avec ce coefficient et empêche cette fuite.

\begin{figure}[htbp]
\centering
\begin{tikzpicture}
\begin{axis}[courseplot,height=4.8cm,axis x line=middle,axis y line=none,
  xlabel={$x$},xmin=-3,xmax=3,ymin=-.25,ymax=.45,ytick=\empty,
  xtick={-3,-2,-1,0,1,2,3},
  legend style={at={(.5,1.02)},anchor=south,legend columns=2,draw=none}]
  \addplot[only marks,mark=*,accent,mark size=2.8pt]
    coordinates {(-2.5,0) (-1.8,0) (-1.2,0) (-.5,0)};
  \addlegendentry{$y=0$}
  \addplot[only marks,mark=triangle*,warning,mark size=3.2pt]
    coordinates {(.4,0) (1.1,0) (1.7,0) (2.5,0)};
  \addlegendentry{$y=1$}
  \draw[ink,very thick] (axis cs:0,-.09)--(axis cs:0,.13)
    node[above=2pt]{seuil $x=0$};
\end{axis}
\end{tikzpicture}
\caption{Séparation complète sur une unique caractéristique $x$. Tous les points
appartiennent au même axe : les cercles ($y=0$) sont à gauche du seuil et les
triangles ($y=1$) à droite. Multiplier le coefficient conserve le seuil mais rend
les probabilités toujours plus extrêmes.}
\label{fig:complete-separation}
\end{figure}

\begin{procedure}{Diagnostiquer l'existence d'une solution finie}
\textbf{Entrées :} la matrice $\X$, les cibles $\y$ et l'historique des
paramètres. \textbf{Sortie :} un diagnostic : ajustement ordinaire, redondance
des colonnes ou séparation complète.
\begin{enumerate}
  \item repérer deux colonnes identiques ou proportionnelles : leurs coefficients
        ne peuvent pas être identifiés séparément ;
  \item vérifier si une frontière classe toutes les observations sans erreur ;
  \item pendant l'ajustement, suivre $\|\thv\|_2$ en plus de $J$ : une norme qui
        croît continuellement tandis que $J$ baisse signale une séparation probable ;
  \item dans ce cas, activer une pénalisation $L^2$ et documenter le diagnostic.
\end{enumerate}
\end{procedure}

\section{Régularisation facultative}

\begin{lecture}
La régularisation compare deux objectifs : bien reproduire les observations et
éviter des coefficients excessivement grands. Le paramètre $\lambda$ fixe le
poids du second objectif. Une valeur nulle conserve le modèle initial ; une
valeur croissante contracte progressivement les coefficients vers zéro.
\end{lecture}

Une pénalisation quadratique donne
\[
  J_\lambda(\thv)=J(\thv)+\frac{\lambda}{2m}\sum_{j=1}^n\theta_j^2,
  \qquad
  \nabla J_\lambda=\nabla J+\frac{\lambda}{m}(0,\theta_1,\ldots,\theta_n)^\top.
\]
Cette forme correspond à la régression logistique pénalisée par une norme
quadratique \cite[sec.~4.4.4]{hastie2009}. Le symbole
\[
  \|\mathbf w\|_2^2=\theta_1^2+\cdots+\theta_n^2
\]
est le carré de la norme euclidienne des coefficients hors biais. Ajouter cette
quantité au coût signifie qu'un ajustement n'est plus jugé uniquement sur son
accord aux données : entre deux solutions de log-pertes voisines, on préfère
celle dont les coefficients sont moins grands.

Le paramètre $\lambda$ règle le compromis. Lorsque $\lambda=0$, on retrouve le
maximum de vraisemblance non pénalisé. Quand $\lambda$ augmente, les coefficients
sont progressivement ramenés vers zéro. À la limite, le modèle dépend de moins en
moins des caractéristiques et tend vers une prédiction principalement gouvernée
par le biais. Le biais n'est conventionnellement pas pénalisé : même si toutes
les caractéristiques deviennent sans effet, le modèle doit encore pouvoir
représenter la fréquence moyenne de la classe positive.

La réduction de variance signifie ici que de petites
modifications de l'échantillon d'apprentissage provoquent généralement de moins
grandes modifications des coefficients et des prédictions. Cette stabilité est
achetée au prix d'un biais statistique supplémentaire : le modèle pénalisé
n'essaie volontairement pas d'ajuster aussi étroitement les données observées.
Le meilleur compromis n'est pas connu à l'avance.

\begin{figure}[htbp]
\centering
\begin{tikzpicture}
\begin{axis}[courseplot,xmode=log,
  xlabel={intensité de régularisation $\lambda$ (échelle logarithmique)},
  ylabel={valeur du coefficient},xmin=.01,xmax=100,ymin=-2.2,ymax=2.8,
  grid=major,grid style={gridgray!40},legend style={at={(.5,1.03)},
  anchor=south,legend columns=3,draw=none},samples=160]
  \addplot[accent,very thick,domain=.01:100]{2.5/(1+x^.55)};
  \addlegendentry{$\theta_1$}
  \addplot[warning,very thick,domain=.01:100]{-1.9/(1+x^.48)};
  \addlegendentry{$\theta_2$}
  \addplot[ink,dashed,very thick,domain=.01:100]{1.25/(1+x^.7)};
  \addlegendentry{$\theta_3$}
  \addplot[gridgray,thin,domain=.01:100]{0};
\end{axis}
\end{tikzpicture}
\caption{Chemins de coefficients sous pénalisation $L^2$ : leur amplitude
décroît lorsque $\lambda$ augmente.}
\label{fig:l2-path}
\end{figure}
\FloatBarrier

On compare donc plusieurs valeurs de $\lambda$ sur un ensemble de validation ou
par validation croisée. Choisir $\lambda$ sur le test reviendrait à adapter le
modèle aux réponses du test ; celui-ci ne fournirait alors plus une évaluation
indépendante. La figure~\ref{fig:lambda-validation} résume cette logique.

\begin{figure}[htbp]
\centering
\begin{tikzpicture}
\begin{axis}[courseplot,xmode=log,
  xlabel={$\lambda$},ylabel={log-perte},xmin=.01,xmax=100,ymin=.2,ymax=1.05,
  grid=major,grid style={gridgray!40},legend pos=north west,samples=180]
  \addplot[accent,very thick,domain=.01:100]{.25+.12*ln(1+x)};
  \addlegendentry{apprentissage}
  \addplot[warning,very thick,domain=.01:100]{.43+.055*(ln(x)-.55)^2};
  \addlegendentry{validation}
  \draw[dashed,ink] (axis cs:1.73,.2)--(axis cs:1.73,.9)
    node[pos=.8,right]{$\lambda$ retenu};
\end{axis}
\end{tikzpicture}
\caption{Sélection de $\lambda$ par minimisation de la log-perte de validation ;
le test final reste exclu de cette sélection.}
\label{fig:lambda-validation}
\end{figure}

\begin{procedure}{Sélectionner la pénalisation}
\textbf{Entrées :} une grille explicite, par exemple
$\lambda\in\{0,10^{-2},10^{-1},1,10\}$, et une partition de validation commune.
\textbf{Sortie :} une valeur de $\lambda$ choisie avant le test.
\begin{enumerate}
  \item réinitialiser et ajuster un modèle pour chaque $\lambda$ sur les mêmes
        lignes d'apprentissage ;
  \item calculer pour chacun la log-perte sur les mêmes lignes de validation ;
  \item construire le tableau $(\lambda,J_{\mathrm{val}})$ et choisir la plus
        petite perte, avec une règle de départage fixée en cas d'égalité ;
  \item enregistrer $\lambda$ et les autres hyperparamètres retenus ;
  \item appliquer ensuite le protocole final et consulter le test une seule fois.
\end{enumerate}
\end{procedure}

\begin{resultat}[Quantités recalculées pendant l'ajustement]
À chaque itération, les paramètres courants produisent les scores, les
probabilités, la log-perte moyenne et le gradient. La log-perte mesure l'état de
l'ajustement ; le gradient produit la correction suivante. La régularisation
ajoute sa contribution au coût et au gradient avant cette correction.
\end{resultat}
